\documentclass[a4paper,12pt,oneside]{article}

\usepackage[T2A]{fontenc}
\usepackage[utf8]{inputenc}
\usepackage[russian]{babel}

\usepackage{lmodern}
\usepackage{microtype}

\usepackage{amsmath,amssymb,amsthm,mathtools}

\usepackage{geometry}
\geometry{top=25mm,left=25mm,right=20mm,bottom=25mm}

\usepackage{graphicx}

\usepackage{enumitem}
\setlist[enumerate]{leftmargin=*, itemsep=0.35\baselineskip, topsep=0.35\baselineskip}
\setlist[itemize]{leftmargin=*, itemsep=0.2\baselineskip, topsep=0.25\baselineskip}

\newlist{alphlist}{enumerate}{1}
\setlist[alphlist]{label=\textbf{\alph*)}, leftmargin=*, align=left}

\newlist{numlist}{enumerate}{1}
\setlist[numlist]{label=\textbf{\arabic*)}, leftmargin=*, align=left}

\newcommand{\Task}[1]{%
  \subsection*{#1}%
  \vspace{-0.2\baselineskip}%
}

\setlength{\parindent}{1.25em}
\setlength{\parskip}{0pt}

\newcommand{\xor}{\oplus}

\begin{document}

\begin{center}
    {\Large\bfseries Полином Жегалкина и суперпозиции}\\[2mm]
\end{center}
\vspace{2mm}
\hrule
\vspace{4mm}

\Task{Задание 1. Метод неопределённых коэффициентов}
Пусть для функции $f(x_1,x_2,\dots,x_n)$ задана таблица истинности. Запишем сначала данную функцию в виде полинома Жегалкина с неопределёнными коэффициентами:
\[
f(x_1,\dots,x_n)=
a_0 \xor a_1x_1 \xor \dots \xor a_nx_n \xor a_{12}x_1x_2 \xor \dots \xor a_{12\dots n}x_1x_2\dots x_n.
\]
Затем по очереди подставляем всевозможные наборы в порядке увеличения количества единиц и находим коэффициенты с учётом того, что $a \xor 1=\overline{a}$, а $a \xor 0=a$. За каждую подстановку находим только один коэффициент.
\[
\includegraphics[width=0.5\linewidth]{tab.jpg}
\]

\Task{Задание 2. Метод треугольника и метод Паскаля}
Построить полином Жегалкина:
\begin{alphlist}
    \item методом треугольника;
    \item методом Паскаля.
\end{alphlist}

\Task{Задание 3. Метод эквивалентных преобразований}
Построить полином Жегалкина для функции методом эквивалентных преобразований:
\[
f(x,y)=x\to\bigl(y\to \overline{x}y\bigr).
\]

\Task{Задание 4. Суперпозиция и упрощение}
Построить и упростить формулу функции $h$, учитывая:
\[
h=f\bigl(x_1,\,g(x_1,x_2),\,x_3\bigr),
\]
где
\[
f(x_1,x_2,x_3)=x_1+x_2+x_2x_3,
\qquad
g(x_1,x_2)=x_1x_2.
\]

\Task{Задание 5. Суперпозиция функций, заданных векторами}
Функции $f(x_1,x_2)$ и $g(x_1,x_2)$ заданы вектором значений, требуется получить суперпозицию $h$.
\begin{numlist}
    \item $f(x_1,x_2)=(1011),\quad g(x_1,x_2)=(0111)$
    \begin{itemize}
        \item $h=f\bigl(x_1,\,g(x_1,x_2)\bigr)$
        \item $h=f\bigl(f(x_1,\,g(x_2,x_2)),\,g(x_1,x_2)\bigr)$
    \end{itemize}

    \item $f(x_1,x_3)=(1010),\quad g(x_2,x_3)=(0110)$
    \begin{itemize}
        \item $h=f\bigl(x_1,\,g(x_2,x_3)\bigr)$
        \item $h=g\bigl(g(x_2,x_3),\,f(x_1,x_1)\bigr)$
    \end{itemize}
\end{numlist}

\end{document}
